% 2.4 — brief en documento/investigacion/2-4-analitica-predictiva.md
\subsection{ANALÍTICA PREDICTIVA}
\subsubsection{DEFINICIÓN Y APLICACIONES}

\textcite{Davenport2007} sostienen que las empresas que compiten con éxito no se limitan a acumular datos, sino que construyen sus estrategias alrededor del conocimiento que extraen de ellos, combinando análisis cuantitativo, modelado predictivo, liderazgo orientado a los datos e infraestructura tecnológica. Interesa que el modelado predictivo figure allí como un componente de la analítica, y no como una disciplina aparte.

Dentro de ese campo, la clasificación más difundida distingue tres niveles. \textcite{Delen2013} separan la analítica descriptiva, que resume lo que ocurrió, de la predictiva, que estima qué podría ocurrir a continuación, y de la prescriptiva, que busca los mejores pasos a seguir. \textcite{Lepenioti2020} confirman en su revisión que el grueso del trabajo académico e industrial se concentra todavía en los dos primeros niveles y presentan a la prescriptiva como el siguiente escalón de madurez analítica, orientado a anticipar decisiones optimizadas. La tríada sirve para ubicar con precisión lo que este proyecto construye: los reportes de Power BI cubren el nivel descriptivo sobre los datos del ERP, los modelos de aprendizaje automático cubren el predictivo, y el nivel prescriptivo, que supondría recomendar acciones correctivas, queda fuera del alcance como trabajo futuro.

Que el objetivo sea predecir tiene consecuencias metodológicas que no siempre se hacen explícitas. \textcite{Shmueli2010} observa que en muchas disciplinas domina el modelado explicativo, orientado a probar relaciones causales, y que con frecuencia se asume que un modelo con alto poder explicativo tiene también alto poder predictivo; la autora muestra que esa confusión es común y que la distinción afecta cada paso del proceso de modelado, desde el diseño y la recolección de datos hasta la selección de variables, la elección del método y la evaluación. Para este trabajo la implicancia es directa: los modelos se juzgan por su error sobre observaciones no vistas, con las métricas y la validación cruzada descritas en la sección~2.2, y no por la significancia o la interpretabilidad de sus coeficientes. Esa prioridad es la que justifica emplear \textit{gradient boosting} aunque sea menos transparente que una regresión lineal.

La analítica predictiva aplicada a operaciones exige, además del método, conocimiento del negocio. \textcite{Waller2013}, al presentarla como una transformación del diseño y la gestión de cadenas de suministro, insisten en que la ciencia de datos requiere tanto dominio del área como un conjunto amplio de habilidades cuantitativas. En el contexto de los sistemas empresariales, \textcite{MLdrivenERP2023} describen la integración de algoritmos de aprendizaje automático en el ERP como el medio para producir predicciones más precisas y decisiones basadas en datos, con el requisito de que los modelos resulten comprensibles para los interesados. El presente proyecto sigue esa línea: el conocimiento del dominio EPC (certificaciones, avance, presupuesto) guía la construcción de variables a partir del histórico del portal anterior y de las planillas de control, y las predicciones se integran al ERP mediante su propio tablero, de modo que llegan al jefe de proyecto en la misma herramienta en que gestiona su trabajo.

\subsubsection{ANALÍTICA PREDICTIVA EN GESTIÓN DE PROYECTOS}

La técnica de pronóstico clásica en control de proyectos es la gestión del valor ganado (\textit{earned value management}, EVM), cuyas magnitudes, variaciones e índices se presentan en la sección~2.5.2. En lo que aquí interesa, el EVM proyecta el resultado final extrapolando el desempeño observado hasta la fecha.

Esa extrapolación tiene límites conocidos. \textcite{Vandevoorde2006} recuerdan que el EVM se desarrolló originalmente para gestionar costos y que su empleo para pronosticar la duración ha sido escaso; comparan los indicadores clásicos de cronograma con los del cronograma ganado (\textit{earned schedule}, ES) y contrastan tres métodos de pronóstico de duración sobre un ejemplo y sobre datos reales. \textcite{Lipke2009} señalan igualmente que el método ofrece desde su origen procedimientos para pronosticar el costo final pero no la duración, e integran el valor ganado, el cronograma ganado y límites de confianza estadísticos para derivar el rango de resultados probables del costo y la duración finales. Para uno de los tres casos de este trabajo, el retraso en la fecha de cierre, ese antecedente es relevante: pronosticar la duración con EVM no es un problema nuevo.

La literatura reciente propone reemplazar la extrapolación de índices por modelos aprendidos del histórico. \textcite{Wauters2016} comparan cinco métodos de inteligencia artificial para predecir la duración final de un proyecto y encuentran que superan a los mejores métodos EVM/ES, sobre todo en las etapas temprana y media, a condición de que los conjuntos de entrenamiento y de prueba sean al menos similares; al aumentar la discrepancia entre ambos, los métodos de IA pierden ventaja. En el frente del costo, \textcite{MLcostForecasting2022} construyen un modelo de red neuronal recurrente (LSTM) alimentado con siete características que incluyen factores de cronograma y desempeño y sus medias móviles, y reportan, sobre 300 experimentos, estimaciones más precisas que el modelo tradicional basado en índices del EVM\@. \textcite{AutoMLpipeline2025} generalizan el enfoque con un pipeline automatizado de balanceo, mejora de características, entrenamiento y evaluación, probado con 30 algoritmos sobre 50 proyectos reales de construcción; sus modelos superan al EVM y a la gestión de cronograma convencionales en exactitud, precisión y oportunidad, y advierten que la literatura previa descuidó el sobreajuste.

Tres consecuencias de esa evidencia orientan el diseño del módulo predictivo. Primera, la condición de similitud señalada por \textcite{Wauters2016} es el argumento técnico para que la capa Silver del pipeline de datos mapee el esquema histórico del portal anterior (MySQL) y el del ERP nuevo (PostgreSQL) a las mismas variables: un modelo entrenado sobre proyectos cerrados solo puede predecir sobre proyectos activos si ambos se describen del mismo modo. Segunda, los indicadores de desempeño y sus medias móviles, empleados como características por \textcite{MLcostForecasting2022}, son candidatos naturales a variables predictoras en la medida en que el histórico permita reconstruirlos, lo que se determina en la exploración de datos del Capítulo III\@. Tercera, la advertencia sobre el sobreajuste de \textcite{AutoMLpipeline2025} refuerza el protocolo de evaluación con conjunto de prueba independiente y validación cruzada ya adoptado en la sección~2.2.

\subsubsection{PREDICCIÓN DE DESVIACIONES Y COSTOS}

Aprender el sobrecosto del histórico supone que tiene un patrón. \textcite{Flyvbjerg2002}, sobre una muestra de 258 proyectos de infraestructura de transporte por 90 mil millones de dólares, concluyen que las estimaciones usadas para decidir su ejecución son sistemáticamente engañosas y que la subestimación no se explica por error, sino por tergiversación estratégica. Aunque se trata de obra pública y no de proyectos EPC, el hallazgo sostiene la premisa de este trabajo: si el sobrecosto tiene un patrón, puede aprenderse del histórico. En el sector EPC, \textcite{EPCleadingIndicators2018} revisan los indicadores líderes de desempeño de costo y cronograma, con poca consistencia entre ellos y una concentración de los estudios en la fase de construcción. La noción de indicador líder es la que el módulo predictivo operacionaliza como variable predictora.

Sobre qué técnica emplear, las revisiones del área convergen en los métodos basados en árboles. \textcite{Elmousalami2020} revisa las técnicas de inteligencia artificial para la estimación paramétrica de costos en construcción, desde la lógica difusa y las redes neuronales hasta los árboles de decisión y el \textit{gradient boosting}, y en un caso de estudio con veinte técnicas obtiene la mejor exactitud con XGBoost (MAPE de 9,09\,\% y $R^2$ ajustado de 0,929), resultado que fue objeto de una discusión posterior en la misma revista. Discute, además, el tratamiento de valores faltantes y atípicos, cuestión pertinente para un histórico cuya calidad de registro está por relevar\@. \textcite{AIcostEstimation2024}, en una revisión sistemática de 39 artículos (2016--2024), encuentran que las redes neuronales son la técnica más frecuente (26,3\,\% de los estudios) y que los árboles de decisión y el \textit{gradient boosting} ofrecen mejoras sustanciales, mientras que la regresión conviene cuando las relaciones son lineales; los rangos de exactitud que reportan agregan métricas heterogéneas, por lo que se toman como tendencia y no como referencia comparable.

Dos trabajos recientes ilustran la arquitectura de solución que este proyecto adopta. \textcite{XGBoostCostOverrun2026} combinan XGBoost con valores SHAP para predecir sobrecostos a partir de variables como tamaño del proyecto, costo estimado, presión de cronograma y cambios de diseño, con validación cruzada de cinco particiones y evaluación con MAE y $R^2$; su conjunto de datos, sin embargo, consta de mil observaciones simuladas, por lo que se cita como ejemplo de método y no como evidencia de desempeño. \textcite{Shen2025} parten de que los modelos basados solo en datos no incorporan el conocimiento experto del dominio y proponen, para predecir retrasos, un marco sobre XGBoost con características derivadas de reglas, penalizaciones en la función de pérdida y reglas de postprocesamiento, que sobre datos reales mejora la exactitud y reduce los falsos negativos. El cuadro~\ref{tab:estudios-predictiva} resume los estudios revisados.

\begin{table}[H]
  \centering
  \caption{Estudios revisados sobre predicción de desviaciones, costos y plazos en proyectos}\label{tab:estudios-predictiva}
  \begin{tabular}{@{}>{\RaggedRight\arraybackslash}p{0.22\textwidth}>{\RaggedRight\arraybackslash}p{0.17\textwidth}>{\RaggedRight\arraybackslash}p{0.13\textwidth}>{\RaggedRight\arraybackslash}p{0.40\textwidth}@{}}
    \toprule
    \textbf{Estudio} & \textbf{Técnica} & \textbf{Variable} & \textbf{Dato relevante} \\
    \midrule
    \textcite{Vandevoorde2006} & Índices EVM y ES & Duración & El EVM se creó para costos y se usó poco para pronosticar duración; compara tres métodos sobre datos reales \\
    \textcite{Wauters2016} & Cinco métodos de IA & Duración & Superan a EVM/ES si entrenamiento y prueba son similares; pierden ventaja al crecer la discrepancia \\
    \textcite{MLcostForecasting2022} & LSTM & Costo & Siete características con medias móviles de indicadores; 300 experimentos; más preciso que el EVM por índices \\
    \textcite{AutoMLpipeline2025} & Pipeline con 30 algoritmos & EAC & 50 proyectos reales de construcción; supera a EVM y gestión de cronograma; alerta sobre sobreajuste \\
    \textcite{EPCleadingIndicators2018} & Revisión de indicadores líderes & Costo y plazo (EPC) & Más de la mitad de los proyectos con retrasos y escaladas de costo importantes \\
    \textcite{Elmousalami2020} & Veinte técnicas de IA & Costo & XGBoost con la mejor exactitud: MAPE 9,09\,\%, $R^2$ ajustado 0,929 \\
    \textcite{AIcostEstimation2024} & Revisión sistemática (39 artículos) & Costo & ANN más frecuente (26,3\,\%); árboles y \textit{gradient boosting} con mejoras sustanciales \\
    \textcite{XGBoostCostOverrun2026} & XGBoost + SHAP & Sobrecosto & CV de 5 particiones, MAE y $R^2$; 1\,000 observaciones simuladas \\
    \textcite{Shen2025} & XGBoost con reglas de dominio & Retraso & Mejora la exactitud y reduce falsos negativos sobre datos reales \\
    \bottomrule
  \end{tabular}
  \fuente{Elaboración propia a partir de las fuentes citadas.}
\end{table}

De esta revisión se desprende el planteo de los tres casos del proyecto: la desviación en certificaciones como problema de clasificación, y el sobrecosto y el retraso en la fecha de cierre como problemas de regresión, los tres con \textit{gradient boosting} sobre el histórico del portal anterior y las planillas de control. Ninguna de las fuentes revisadas aporta cifras de sobrecosto para el sector de petróleo, gas y minería en Bolivia, por lo que la magnitud de las desviaciones en ISI Mustang solo puede establecerse con sus propios datos. Se prevé incorporar reglas de negocio como variables y validaciones de consistencia, como en \textcite{Shen2025}. Las predicciones se muestran únicamente en el tablero del ERP; los reportes de Power BI se mantienen en el nivel descriptivo.
